\documentclass[a4paper]{article}
\usepackage{ctex}
\ctexset{proofname=\heiti{证明}}
\usepackage{amsmath,amssymb,amsthm}
\usepackage{mathtools}
\usepackage{bm}
\usepackage{enumitem}
\usepackage{booktabs}
\usepackage{array}
\usepackage{graphicx}

\IfFileExists{iidef.sty}{
    \usepackage[thehwcnt = 2]{iidef}
    \thecourseinstitute{清华大学}
    \thecoursename{人工智能原理}
    \theterm{2026年春季学期}
    \hwname{作业}
    \slname{\heiti{解}}
}{
    \makeatletter
    \newcommand{\thecourseinstitute}[1]{\def\@courseinstitute{#1}}
    \newcommand{\thecoursename}[1]{\def\@coursename{#1}}
    \newcommand{\theterm}[1]{\def\@term{#1}}
    \newcommand{\hwname}[1]{\def\@hwname{#1}}
    \newcommand{\slname}[1]{\def\@slname{#1}}
    \thecourseinstitute{清华大学}
    \thecoursename{人工智能原理}
    \theterm{2026年春季学期}
    \hwname{作业}
    \slname{\heiti{解}}
    \newcommand{\courseheader}{
        \begin{center}
            {\Large \heiti \@courseinstitute \quad \@coursename}\\[0.5em]
            {\large \@term \quad \@hwname 2 \quad \@slname}
        \end{center}
        \vspace{1em}
    }
    \newcommand{\name}[1]{\noindent #1 \par\vspace{1em}}
    \newenvironment{solution}{\par\noindent\textbf{解：}\par}{\par}
    \makeatother
}

\begin{document}
\courseheader
\name{姓名：杜一郎\qquad 学号：2023011618}

\begin{enumerate}
\setlength{\itemsep}{3\parskip}

\item Minimax 与 alpha-beta 剪枝

\begin{solution}
\textbf{(1) Minimax}

底层 6 个 MAX 节点从左到右记为 \(M_1,\ldots,M_6\)。它们的取值为
\[
\begin{array}{c|c|c}
\toprule
\text{节点} & \text{子 MIN 值} & \text{MAX 值}\\
\midrule
M_1 & \min(5,8)=5,\ \min(2,9)=2 & 5\\
M_2 & 7,\ 0 & 7\\
M_3 & 4,\ \min(6,1)=1 & 4\\
M_4 & \min(3,6)=3,\ 8 & 8\\
M_5 & 9,\ 7 & 9\\
M_6 & 2,\ \min(1,5)=1 & 2\\
\bottomrule
\end{array}
\]
左侧 MIN 节点的值为
\[
\min(5,7,4)=4,
\]
右侧 MIN 节点的值为
\[
\min(8,9,2)=2.
\]
根节点是 MAX，因此
\[
\max(4,2)=4.
\]
所以根节点最终值为 \(4\)，应选择左分支。

\textbf{(2) alpha-beta 剪枝}

按从左到右的顺序生成子节点：
\begin{itemize}
    \item 左侧 MIN 的第二个 MAX 节点中，先得到值 \(7\)，已经不可能让父 MIN 优于当前 \(\beta=5\)，剪去叶子 \(0\)。
    \item 右侧 MIN 的第一个 MAX 节点中，子 MIN 看到叶子 \(3\) 后已有 \(\beta\le\alpha\)，剪去叶子 \(6\)。
    \item 右侧 MIN 的第二个 MAX 节点中，先得到 \(9\)，已经有 \(\alpha\ge\beta\)，剪去叶子 \(7\)。
    \item 右侧 MIN 的第三个 MAX 节点中，子 MIN 看到叶子 \(1\) 后满足 \(\beta\le\alpha\)，剪去叶子 \(5\)。
\end{itemize}
被剪掉的叶节点共 \(4\) 个。alpha-beta 剪枝不改变 minimax 值，根节点仍为 \(4\)，仍选择左分支。
\end{solution}

\item MCTS

\begin{solution}
UCB1 公式为
\[
U_{\mathrm{black}}(i)=\frac{w_i}{n_i}+c\sqrt{\frac{\ln N}{n_i}},
\]
\[
U_{\mathrm{white}}(i)=1-\frac{w_i}{n_i}+c\sqrt{\frac{\ln N}{n_i}}.
\]

\textbf{(1) \(c=0.2\)}

在根节点 \(s_0\)，\(N=40\)，黑方选择：
\[
\begin{array}{c|c}
\toprule
\text{子节点} & \text{UCB}\\
\midrule
s_1 & 0.917\\
s_2 & 0.859\\
s_3 & 0.511\\
\bottomrule
\end{array}
\]
所以先选 \(s_1\)。在 \(s_1\) 处轮到白方选择，\(N=25\)：
\[
\begin{array}{c|c}
\toprule
\text{子节点} & \text{UCB}\\
\midrule
s_4 & 0.541\\
s_5 & 0.564\\
s_6 & 0.093\\
\bottomrule
\end{array}
\]
因此选择 \(s_5\)，Selection 路径为
\[
s_0\rightarrow s_1\rightarrow s_5.
\]
扩展 \(e_1\)，若模拟结果为黑方获胜，则回传后
\[
s_0:29/40\rightarrow 30/41,
\]
\[
s_1:21/25\rightarrow 22/26,
\]
\[
s_5:4/7\rightarrow 5/8,
\]
\[
e_1:1/1.
\]

\textbf{(2) \(c=1.2\)}

在根节点 \(s_0\)：
\[
\begin{array}{c|c}
\toprule
\text{子节点} & \text{UCB}\\
\midrule
s_1 & 1.301\\
s_2 & 1.585\\
s_3 & 1.190\\
\bottomrule
\end{array}
\]
所以选择 \(s_2\)。在 \(s_2\) 处轮到白方选择，\(N=7\)：
\[
\begin{array}{c|c}
\toprule
\text{子节点} & \text{UCB}\\
\midrule
s_7 & 1.674\\
s_8 & 1.087\\
s_9 & 1.684\\
\bottomrule
\end{array}
\]
因此选择 \(s_9\)，Selection 路径为
\[
s_0\rightarrow s_2\rightarrow s_9.
\]
扩展 \(e_2\)，若模拟结果为黑方获胜，则回传后
\[
s_0:29/40\rightarrow 30/41,
\]
\[
s_2:5/7\rightarrow 6/8,
\]
\[
s_9:1/2\rightarrow 2/3,
\]
\[
e_2:1/1.
\]

\textbf{(3) 探索系数比较}

\(c=0.2\) 时探索项较小，选择更偏向当前胜率高的节点，所以走到
\[
s_0\rightarrow s_1\rightarrow s_5.
\]
\(c=1.2\) 时探索项较大，访问次数少的节点会得到更高加成，所以走到
\[
s_0\rightarrow s_2\rightarrow s_9.
\]
因此 \(c\) 越大，MCTS 在 Selection 阶段越偏探索；\(c\) 越小，越偏利用当前估计胜率。
\end{solution}

\end{enumerate}
\end{document}
